\section*{Chapter \ref{ch:iv1:stochastic-calculus} --- Stochastic Calculus}
\markboth{\protect\omTocMark{}{Answers: Stochastic Calculus}}{\protect\omTocMark{}{Answers: Stochastic Calculus}}

\begin{solution}{iq:iv1:stochastic-calculus:1}
$\operatorname{Cov}(W_s, W_t) = \min(s, t)$. The variance is 5, the correlation $\frac12$, and the probability $\frac38$.

For $s \le t$, write $W_t = W_s + (W_t - W_s)$, whose second term is independent of $W_s$. That leaves $\operatorname{Var}W_s = s$.

Then $\operatorname{Var}(W_1 + W_2) = 1 + 2 + 2(1) = 5$. The correlation of $W_1$ and $W_4$ is $1/\sqrt{1 \times 4} = \frac12$.

For two jointly normal variables with mean zero and correlation $r$, the chance that both are positive is $\frac14 + \arcsin(r)/(2\pi)$\index{orthant probability}. Here $r = 1/\sqrt{2}$, so $\arcsin r = \pi/4$, and the probability is $\frac14 + \frac18$. A simulation gives $0.375$.

\iqwrong{Treating $W_1$ and $W_2$ as independent, which gives a variance of 3 and a probability of $\frac14$. They overlap on $[0, 1]$.}

\iqfollow{What is the variance of $W_1 + \dots + W_{10}$?}{$\sum_{i, j}\min(i, j) = 1^2 + 2^2 + \dots + 10^2 = 385$, a useful identity.}

\iqwhy{The covariance formula is the source of everything else, and the question checks that it can be used and not only recited.}
\end{solution}

\begin{solution}{iq:iv1:stochastic-calculus:2}
Martingales: $W_t$, $W_t^2 - t$, $e^{W_t - t/2}$ and $W_t^3 - 3tW_t$. Not martingales: $W_t^2$, $tW_t$ and $W_t^3$. The test is whether $f_t + \frac12f_{xx} = 0$.

By Ito's lemma\index{Ito's lemma}, $df(t, W_t) = (f_t + \frac12f_{xx})\,dt + f_x\,dW$. A martingale needs the $dt$ term to vanish.

$W_t$ gives $0$. $W_t^2$ gives $1$, so it has a drift, and subtracting $t$ removes it. $tW_t$ gives $W_t$, which is not zero. $e^{W_t - t/2}$ gives $-\frac12 f + \frac12f = 0$. $W_t^3$ gives $3W_t$, and $W_t^3 - 3tW_t$ gives $-3W_t + 3W_t = 0$. Gaussian quadrature of the conditional expectations confirms the seven verdicts.

A rule of thumb for $tW_t$: its mean is 0, which may suggest a martingale. Constant mean is necessary, and not sufficient. Its conditional mean given time $s$ is $tW_s$, not $sW_s$.

\iqwrong{Saying that $W_t^3$ is a martingale because its mean is always zero. Given $W_s$, the mean of $W_t^3$ is $W_s^3 + 3(t - s)W_s$.}

\iqfollow{Find the next polynomial in the family.}{$W_t^4 - 6tW_t^2 + 3t^2$. These are Hermite polynomials in $W_t/\sqrt{t}$, scaled by powers of $\sqrt{t}$.}

\iqwhy{The one-line test replaces seven separate calculations. Knowing it marks a working knowledge of Ito's lemma.}
\end{solution}

\begin{solution}{iq:iv1:stochastic-calculus:3}
$d\ln S = (\mu - \frac12\sigma^2)\,dt + \sigma\,dW$. The mean is $e^{0.1} = 1.105$, the median $e^{0.08} = 1.083$, and the probability $0.46$.

Ito's lemma gives $d\ln S = dS/S - \frac12(dS)^2/S^2$, and $(dS)^2 = \sigma^2S^2\,dt$. So the log has drift $\mu - \frac12\sigma^2 = 8\%$\index{geometric Brownian motion}.

$\ln S_1$ is normal with mean $0.08$ and standard deviation $0.2$. The median is at the mean of the log, so $e^{0.08}$. For the mean of $S$, $\mathbb{E}e^{0.08 + 0.2Z} = e^{0.08 + 0.02} = e^{0.1}$, as the equation for $S$ also shows.

The chance of ending above the mean needs $0.08 + 0.2Z > 0.1$, that is, $Z > 0.1$, with probability $1 - \Phi(0.1) = 0.46$. A simulation of the discretised equation gives these numbers.

With no volatility, mean and median agree. The gap grows with $\sigma^2$.

\iqwrong{Dropping the $-\frac12\sigma^2$, which is the second-order Ito term, and giving a median of $e^{0.1}$.}

\iqfollow{At what volatility is the median growth zero, with $\mu = 10\%$?}{$\sigma^2/2 = 0.1$, so $\sigma = 44.7\%$. The average path still grows at 10\% while the typical path goes nowhere.}

\iqwhy{This is the first thing done with Ito's lemma in finance, and the difference between mean and median is its most useful result.}
\end{solution}

\begin{solution}{iq:iv1:stochastic-calculus:4}
$\mathbb{E}W_t^4 = 3t^2$, $\mathbb{E}e^{aW_t} = e^{a^2t/2}$, $\mathbb{E}e^{2W_1} = e^2 \approx 7.39$, and the correlation is $\frac14$.

Since $W_t = \sqrt{t}Z$, the fourth moment is $t^2\mathbb{E}Z^4 = 3t^2$. Completing the square in the normal integral gives the exponential formula\index{moment generating function}.

For the squares, write $W_4 = W_1 + V$, where $V$ is independent of $W_1$ with variance 3. Then
\[
\mathbb{E}W_1^2W_4^2 = \mathbb{E}W_1^4 + \mathbb{E}W_1^2\,\mathbb{E}V^2 = 3 + 3 = 6.
\]
So the covariance is $6 - 1 \times 4 = 2$. The variances are $2$ and $2 \times 16 = 32$, which gives $2/\sqrt{64} = \frac14$. This is the square of the correlation $\frac12$ of $W_1$ and $W_4$, as is always true of two jointly normal variables with mean zero. A simulation agrees.

\iqwrong{Giving $\mathbb{E}W_t^4 = t^2$, which ignores the kurtosis. Or giving the correlation as $\frac12$, the value for $W_1$ and $W_4$ themselves.}

\iqfollow{What is $\mathbb{E}W_t^6$?}{$15t^3$. The even normal moments are $1, 3, 15, 105$, each the previous times the next odd number.}

\iqwhy{These moments come up in every variance calculation. Quick recall saves effort, and the correlation of squares is a small test of care.}
\end{solution}

\enlargethispage{-3\baselineskip}%
\begin{solution}{iq:iv1:stochastic-calculus:5}
$P(W_4 > 2) = P(W_1 > 1) = 0.1587$, and the annual volatility is $15.9\%$.

$X_t = cW_{t/c^2}$ is continuous, starts at 0, and has independent normal increments with mean zero. Its covariance is $c^2\min(s/c^2, t/c^2) = \min(s, t)$\index{Brownian scaling}. So it is a Brownian motion.

With $c = 2$: $W_4 = 2W'_1$ for a Brownian motion $W'$, so $P(W_4 > 2) = P(W'_1 > 1)$, which is $1 - \Phi(1) = 0.1587$.

Independent returns add their variances, so standard deviations scale with the square root of time: $1\% \times \sqrt{252} = 15.9\%$. If returns are not independent, this figure is wrong, which is why some measure of autocorrelation has to accompany it.

\iqwrong{Multiplying by 252, which gives 252\%. Variance scales with time, and volatility with its root.}

\iqfollow{What is the annual volatility for a monthly volatility of 5\%?}{$5\% \times \sqrt{12} = 17.3\%$.}

\iqwhy{Scaling is the Brownian fact used most on any desk, and proving it with the covariance formula takes one line.}
\end{solution}

\begin{solution}{iq:iv1:stochastic-calculus:6}
The sum of squares has mean 1 and standard deviation $0.045$. The sum of absolute values has a mean of $25.2$. As the steps shrink, the first converges to 1, and the second diverges.

Each increment is normal with variance $1/n$. Its square has mean $1/n$ and variance $2/n^2$. Summing over $n$ independent steps gives mean 1 and variance $2/n$\index{quadratic variation}, so the sum converges to 1 and its standard deviation at $n = 1000$ is $\sqrt{0.002}$.

The absolute value has mean $\sqrt{2/\pi}/\sqrt{n}$, and so the sum has mean $\sqrt{2n/\pi}$, which is $25.2$ here. It grows like $\sqrt{n}$\index{total variation}. So a Brownian path has infinite length but finite quadratic variation, and this is the reason that the extra term in Ito's lemma exists: $(dW)^2$ behaves like $dt$.

A simulation with $20\,000$ paths confirms all three numbers.

\iqwrong{Thinking that the sum of squares is random in the limit, or that it tends to 0 as for a smooth path. For a smooth path the sum of squares goes to 0, while here it stays at $T$, with no randomness left.}

\iqfollow{With 252 daily returns, how noisy is a realised variance?}{The relative standard deviation is $\sqrt{2/252} = 8.9\%$ for constant volatility and normal returns, and more with fat tails.}

\iqwhy{This is the heart of the subject. It explains $(dW)^2 = dt$ with a calculation, and it is the theory behind realised variance.}
\end{solution}

\begin{solution}{iq:iv1:stochastic-calculus:7}
$\int_0^TW\,dW = \frac12(W_T^2 - T)$, with mean 0 and variance $T^2/2$. The left sum converges to this, the right sum to this plus $T$, and the midpoint to $\frac12W_T^2$.

Ito's lemma on $W^2$ gives $d(W^2) = 2W\,dW + dt$, which rearranges to the result\index{Ito integral}. The mean is 0, since $\mathbb{E}W_T^2 = T$. The variance is $\frac14\operatorname{Var}W_T^2 = \frac14 \times 2T^2$, or by the Ito isometry\index{Ito isometry}, $\int_0^T\mathbb{E}W_t^2\,dt = T^2/2$.

The left-end sum is the definition, and its terms have mean zero, as the integrand is known before the increment.

The right sum minus the left sum is $\sum(\Delta W)^2$, which converges to $T$, so the right sum has mean $T$. In the limit, it is $\frac12(W_T^2 + T)$.

The average of the two sums is the midpoint form, and it telescopes exactly: $\sum\frac12(W_i + W_{i + 1})(W_{i + 1} - W_i) = \frac12W_T^2$. This is the Stratonovich integral\index{Stratonovich integral}, which keeps the rules of ordinary calculus, and evaluating $W$ at the true midpoint of each step has the same limit. A run with 1000 steps shows the left, right and averaged sums.

\iqwrong{Writing $\frac12W_T^2$ from ordinary calculus. Its mean is $T/2$, which is impossible for a sum of terms with mean zero.}

\iqfollow{Why does finance use the left endpoint?}{Because a position must be chosen before the price move it earns. A right-end rule would trade on the future, and it would earn exactly the extra $T$.}

\iqwhy{This is the shortest illustration of why stochastic integrals differ from ordinary ones, and it shows that the choice of endpoint has real content.}
\end{solution}

\enlargethispage{-1\baselineskip}%
\begin{solution}{iq:iv1:stochastic-calculus:8}
$I_T$ is normal with mean 0 and variance $T^3/3$, which is 9 at $T = 3$. Its correlation with $W_T$ is $\sqrt{3}/2 = 0.866$.

Integrating by parts\index{integrated Brownian motion}, $d(tW_t) = W_t\,dt + t\,dW_t$, so
\[
I_T = TW_T - \int_0^Tt\,dW_t = \int_0^T(T - t)\,dW_t.
\]
This is an integral of a deterministic function against $dW$, and it is therefore normal with mean 0. The variance by the isometry is $\int_0^T(T - t)^2\,dt = T^3/3$.

As a check, the double integral gives $\int_0^T\int_0^T\min(s, t)\,ds\,dt = T^3/3$, and at $T = 3$ a quadrature also gives 9.

The covariance with $W_T = \int_0^T dW$ is $\int_0^T(T - t)\,dt = T^2/2$. So the correlation is
\[
\frac{T^2/2}{\sqrt{T^3/3 \cdot T}} = \frac{\sqrt{3}}{2},
\]
whatever $T$ is. A simulation with 2000 steps agrees.

\iqwrong{Giving the variance as $T^2/2$ or $T^3$. A units check helps: $I_T$ has the units of $W$ times time, so its variance scales as $T \times T^2$.}

\iqfollow{How does the volatility of a price averaged over a period compare with that of the price itself?}{The average $I_T/T$ has variance $T/3$, so its volatility is the price's divided by $\sqrt{3}$. A 20\% volatility becomes $11.5\%$, which is why Asian options are cheaper.}

\iqwhy{Integrated Brownian motion is one of the standard calculations, and the integration by parts is more insightful than the double integral.}
\end{solution}

\begin{solution}{iq:iv1:stochastic-calculus:9}
$P(M_T \ge a) = 2P(W_T \ge a)$, so $M_T$ has the law of $|W_T|$. $P(M_1 \ge 1) = 0.317$, and $\mathbb{E}M_1 = 0.798$. Checking only at daily closes moves the barrier out by about $0.5826\sigma\sqrt{\Delta}$.

Once $W$ first reaches $a$, the rest of the path is symmetric about $a$\index{reflection principle}, so the end is as likely to lie above $a$ as below it. So $P(M_T \ge a, W_T < a) = P(M_T \ge a, W_T > a) = P(W_T > a)$. Adding gives the factor 2.

So $P(M_T \ge a) = 2(1 - \Phi(a/\sqrt{T}))$, which is the tail of $|W_T|$. At $a = T = 1$ it is $0.3173$. The mean is $\mathbb{E}|W_1| = \sqrt{2/\pi}$.

A path can cross between two closes and return. So the discrete maximum is lower, and the hit probability is lower. The usual correction\index{discrete monitoring} shifts the barrier by $0.5826$ standard deviations of one step, here $0.5826/\sqrt{252} = 0.037$. It gives $2(1 - \Phi(1.037)) = 0.300$. A simulation of daily closes with $100\,000$ paths gives this number, not $0.317$.

\iqwrong{Forgetting the factor 2, or using the continuous formula for a barrier checked once a day. On a one-year barrier, that is about two points of probability.}

\iqfollow{Why does the factor 2 not hold with drift?}{The reflected path is no longer equally likely. With drift, the formula gains a second term with a factor of $e^{2\mu a/\sigma^2}$.}

\iqwhy{This is the standard result for barrier options, and the correction for discrete monitoring is where the theory meets a real term sheet.}
\end{solution}

\begin{solution}{iq:iv1:stochastic-calculus:10}
From the distribution $P(T_a \le t) = 2(1 - \Phi(a/\sqrt{t}))$, the median for $a = 1$ is $2.20$, and $P(T_1 > 100) = 0.080$. The mean is infinite because the tail is of order $1/\sqrt{t}$.

$T_a \le t$ is the event $M_t \ge a$\index{first passage time}, so the reflection principle gives the formula. As $t \to \infty$ it tends to $2(1 - \Phi(0)) = 1$, so the time is finite with probability one.

For the median, $2(1 - \Phi(1/\sqrt{t})) = \frac12$ gives $1/\sqrt{t} = \Phi^{-1}(0.75) = 0.6745$, so $t = 2.198$.

Also, $P(T_1 > 100) = 1 - 2(1 - \Phi(0.1)) = 0.0797$.

For large $t$, $P(T_a > t) \approx 2\phi(0)a/\sqrt{t} = 0.798a/\sqrt{t}$. The mean is the integral of this tail, which diverges. Equivalently, $T_a$ has the law of $a^2/Z^2$. Samples of this form confirm both the median and the tail, and truncated means keep growing as the cutoff rises. A direct run of the path gives a hit by $t = 2.198$ about half the time.

\iqwrong{Assuming that a certain event must have a finite mean waiting time. Here half the paths arrive by $2.2$, while 8\% are still out at 100.}

\iqfollow{How does the median change for $a = 10$?}{By scaling, it is multiplied by $a^2$, so the median becomes about 220.}

\iqwhy{Certain arrival with infinite mean is the key surprise of the subject, and it reappears wherever someone assumes that a mean-zero bet will come back within a known time.}
\end{solution}

\begin{solution}{iq:iv1:stochastic-calculus:11}
The probability of ending at 2 is $\frac13$, and the expected time is 2. In general, these are $a/(a + b)$ and $ab$.

Let $\tau$ be the stopping time. It has finite mean, and the stopped $W$ is bounded, so optional stopping\index{optional stopping} applies.

With $W$: $\mathbb{E}W_\tau = 0$. If $q$ is the probability of ending at $b$, then $qb - (1 - q)a = 0$, so $q = a/(a + b)$. With $a = 1$, $b = 2$, this is $\frac13$.

With $W^2 - t$: $\mathbb{E}W_\tau^2 = \mathbb{E}\tau$. So
\[
\mathbb{E}\tau = qb^2 + (1 - q)a^2 = \frac{ab^2 + ba^2}{a + b} = ab = 2.
\]
A symmetric lattice walk has these same values exactly, and a run with $20\,000$ paths on a lattice of $0.1$ gives $0.33$ and $2.0$.

As $a \to \infty$, $q \to 1$ and $\mathbb{E}\tau = ab \to \infty$: reaching $b$ is certain and takes infinite expected time, which agrees with the one-sided hitting time.

\iqwrong{Using optional stopping with no check. On the one-sided problem it would give $\mathbb{E}W_\tau = 0$, while $W_\tau = b$. It is the bound on the stopped process that makes it valid here.}

\iqfollow{What are the answers for a symmetric interval $\pm a$?}{$\frac12$ and $a^2$. So the typical time to move a distance $a$ is $a^2$, the same scale as the variance.}

\iqwhy{Two martingales and one line each. It is the canonical use of optional stopping, and interviewers listen for the check on when it applies.}
\end{solution}

\begin{solution}{iq:iv1:stochastic-calculus:12}
$e^{-2\mu X_t/\sigma^2}$ is a martingale, which gives $P(\text{ever at } {-a}) = e^{-2\mu a/\sigma^2}$. Here $e^{-1} = 0.368$.

Try $e^{\theta X_t}$. Its drift has the factor $\theta\mu + \frac12\theta^2\sigma^2$, which vanishes at $\theta = -2\mu/\sigma^2$\index{exponential martingale}.

Stop at $\tau$, the first time at $-a$, or at some high level $b$. The martingale is bounded on the interval, so optional stopping gives
\[
1 = qe^{2\mu a/\sigma^2} + (1 - q)e^{-2\mu b/\sigma^2},
\]
where $q$ is the chance of hitting $-a$ first. As $b \to \infty$, the second term vanishes, so $q = e^{-2\mu a/\sigma^2}$.

With $\mu = 0.1$, $\sigma = 0.2$ and $a = 0.2$: $2 \times 0.1 \times 0.2/0.04 = 1$, so the answer is $e^{-1}$. A walk with matching mean and variance has the exact ruin formula, and it gives $0.368$ for small steps. A simulation in steps of $0.01$ years agrees, after the correction for checking at discrete times.

Note the limits: $\mu \to 0$ gives 1, as the driftless path reaches every level, and a large $\mu$ gives a small number. The key ratio is $\mu a/\sigma^2$, so doubling the volatility divides the exponent by four.

\iqwrong{Using a normal tail at one horizon, such as a one-year chance of $-20\%$. That misses all the paths that dip and recover, and it is much smaller than the true figure.}

\iqfollow{What is the chance of ever losing 50\%, with the same numbers?}{$e^{-2.5} = 0.082$. The probability falls exponentially in the size of the loss.}

\iqwhy{This is a practical risk-of-ruin formula. It is derived in three lines with a martingale, and it is often needed in a discussion of sizing.}
\end{solution}

\begin{solution}{iq:iv1:stochastic-calculus:13}
$\mathbb{E}S_T^2 = S_0^2e^{(2\mu + \sigma^2)T}$, the martingale power is $p = 1 - 2\mu/\sigma^2 = -4$, and the probability is $15/15.94 = 0.941$.

Ito's lemma gives
\[
d(S^p) = S^p\left(\left(p\mu + \tfrac12p(p - 1)\sigma^2\right)dt + p\sigma\,dW\right).
\]
So $\mathbb{E}S_T^p = S_0^p\exp((p\mu + \frac12p(p - 1)\sigma^2)T)$, and $p = 2$ gives the second moment, $e^{0.24} = 1.271$ per unit of initial square at $T = 1$.

The drift is zero when $\mu + \frac12(p - 1)\sigma^2 = 0$, or at $p = 0$, which is trivial. So $p = 1 - 2\mu/\sigma^2 = 1 - 5 = -4$\index{scale function}.

Stopping $S^{-4}$ at $2$ or $\frac12$, where it is bounded: $q2^{-4} + (1 - q)2^4 = 1$. So $q = 15/(16 - \frac{1}{16}) = 0.941$.

A second route works with $\ln S$, with drift $0.08$ and volatility $0.2$, exiting at $\pm\ln 2$. The exponential martingale $e^{\theta X}$ with $\theta = -4$ gives the same equation, which is no surprise, since $S^{-4} = e^{-4\ln S}$. A simulation with fine steps returns $0.94$.

\iqwrong{Saying that the probability is $\frac12$, since both levels are a factor of 2 away, or that it follows from $\mu$ alone. The log drift of $0.08$ against a volatility of $0.2$ is what makes this so lopsided.}

\iqfollow{What is $p$ when $\mu = \frac12\sigma^2$?}{$p = 0$, and the martingale is then the log itself, since the log has no drift. The exit probability then is exactly $\frac12$.}

\iqwhy{Finding the power that removes the drift is the fastest way to solve exit problems for geometric processes.}
\end{solution}

\begin{solution}{iq:iv1:stochastic-calculus:14}
The mean is $sx$, and the variance $s(1 - s)$. Given $W_1 = 0$, the probability of reaching 1 is $e^{-2} = 0.135$. In general, for $a \ge \max(x, 0)$, it is $e^{-2a(a - x)}$.

$(W_s, W_1)$ is jointly normal, with covariance $s$ and $\operatorname{Var}W_1 = 1$. The regression gives the conditional mean $sx$ and the variance $s - s^2$\index{Brownian bridge}. So the bridge is pinned at both ends, with the largest uncertainty, $\frac14$, at the midpoint.

By reflection at $a$, a path that ends at $x < a$ and has touched $a$ corresponds to a path that ends at $2a - x$. So the density of paths with both properties is $\phi(2a - x)$, and the conditional probability is the ratio of densities:
\[
\frac{\phi(2a - x)}{\phi(x)} = \exp\left(-\frac{(2a - x)^2 - x^2}{2}\right) = e^{-2a(a - x)}.
\]
With $x = 0$ and $a = 1$, this is $e^{-2}$. Note the limit $x \to a$, where the formula gives 1, as it must. A simulated bridge with 2000 steps gives $0.135$ after the correction for discrete checking.

\iqwrong{Using $2(1 - \Phi(1))$ from the unconditional maximum. That averages over all endpoints. Pinning the end at 0 reduces the chance by more than half.}

\iqfollow{A price moves from 100 to 101 in a day with a daily volatility of 1. What is the chance that it touched 102 between the two closes?}{Here $a = 2$ and $x = 1$, giving $e^{-4} = 0.018$. This is how barrier monitoring is corrected between observation dates.}

\iqwhy{This formula is in every barrier pricer and in path-dependent simulation. The derivation by reflection is short enough for an interview.}
\end{solution}

\begin{solution}{iq:iv1:stochastic-calculus:15}
$X_t = X_0e^{-\kappa t} + \sigma\int_0^te^{-\kappa(t - s)}\,dW_s$. The mean is $0.406$, the variance $0.245$, the stationary variance $\frac14$, and the half-life $0.347$.

Multiply by the integrating factor\index{Ornstein-Uhlenbeck process}: $d(e^{\kappa t}X) = e^{\kappa t}\sigma\,dW$. Integrating gives the solution.

The mean is $X_0e^{-\kappa t}$, here $3e^{-2} = 0.406$. By the isometry, the variance is
\[
\sigma^2\int_0^te^{-2\kappa(t - s)}\,ds = \frac{\sigma^2}{2\kappa}\left(1 - e^{-2\kappa t}\right),
\]
which is $\frac14(1 - e^{-4}) = 0.245$ at these values. As $t$ grows, this tends to the stationary value $\sigma^2/(2\kappa) = \frac14$. The half-life solves $e^{-\kappa t} = \frac12$, so $t = \ln 2/\kappa$.

Limits: for small $t$, the variance is close to $\sigma^2t$, as for plain Brownian motion. As $\kappa \to 0$, the same is true for all $t$. Both an exact AR(1) sampling and an Euler run give $0.406$ and $0.245$.

\iqwrong{Forgetting that the mean-reversion also reduces the noise's contribution, and giving the variance as $\sigma^2t$. It is bounded, which is the point of the model.}

\iqfollow{A daily series has AR(1) coefficient $0.98$. What are $\kappa$ and the half-life?}{$\kappa = -\ln 0.98 = 0.0202$ per day, and the half-life is $34.3$ days. This is the standard calculation for a mean-reverting spread.}

\iqwhy{This is the second-most used stochastic process in finance, the base of many pairs and rates models.}
\end{solution}

\begin{solution}{iq:iv1:stochastic-calculus:16}
$W^2 = \rho W^1 + \sqrt{1 - \rho^2}\,W^\perp$, $\mathbb{E}W^1_tW^2_t = \rho t$, and the two values are 3 and $1.5$.

The construction has unit variance, and the correlation with $W^1$ is $\rho$\index{correlated Brownian motions}.

Ito's lemma gives $d(W^1W^2) = W^1\,dW^2 + W^2\,dW^1 + \rho\,dt$, since $dW^1dW^2 = \rho\,dt$. The stochastic integrals have mean zero, so the mean of the product is $\rho t$.

Then $\operatorname{Var}(W^1 + W^2) = t + t + 2\rho t = 3$.

For the last one, $\mathbb{E}[X^2Y^2] = \sigma_X^2\sigma_Y^2 + 2\operatorname{Cov}(X, Y)^2$ for jointly normal variables with mean zero, by Isserlis' theorem\index{Isserlis' theorem}. That gives $t^2(1 + 2\rho^2) = 1.5$. Direct substitution checks this: writing $W^2 = \frac12W^1 + \sqrt{\frac34}V$ with $V$ independent gives $\frac14 \times 3 + \frac34 \times 1 = 1.5$. A simulation confirms this. At $\rho = 0$, the result is 1, the product of the variances, and at $\rho = 1$ it is 3, the fourth moment.

\iqwrong{Multiplying the variances to get 1, which ignores the correlation of squares. Or writing $W^2 = \rho W^1 + (1 - \rho)W^\perp$, which does not have unit variance.}

\iqfollow{What is the variance of the sum at $\rho = -\frac12$?}{$2t(1 + \rho) = 1$ at $t = 1$. That is the basic diversification formula.}

\iqwhy{Every multi-asset simulation is built this way. The Ito cross-term is where people go wrong.}
\end{solution}

\begin{solution}{iq:iv1:stochastic-calculus:17}
The identity follows from the shift of the density. The tilted estimator reduces the relative standard error from 18\% to about $0.2\%$, a factor of more than 60.

With $Z = W_T$ and density $\phi_T$,
\[
\mathbb{E}[f(W_T + \mu T)] = \int f(x + \mu T)\phi_T(x)\,dx = \int f(y)\phi_T(y - \mu T)\,dy,
\]
and completing the square gives $\phi_T(y - \mu T) = \phi_T(y)e^{\mu y - \mu^2T/2}$\index{Girsanov's theorem}. So the drift can be moved into a weight. That is Girsanov's theorem for a single time, and it extends to paths.

As a check, with $\mu = T = 1$ and $f = \mathbf{1}(x > 0)$, both forms give $\Phi(1) = 0.8413$.

For $P(W_1 > 4) = 3.17 \times 10^{-5}$, a plain simulation with $10^6$ paths gets about 32 hits, and so has a relative error of $18\%$.

Reverse the identity: simulate $Y = Z + 4$ and average $\mathbf{1}(Y > 4)e^{-4Y + 8}$\index{importance sampling}. Now half the paths land in the event, and each carries a small weight. On a run, this reduces the relative error to about $0.2\%$, matching an integral for the second moment. Plain sampling would need several thousand times as many paths for the same accuracy.

Too large a shift hurts again, because the weights become erratic. For a tail event, the good choice puts the mean at the boundary of the event.

\iqwrong{Dropping the weight, or using $e^{4Y - 8}$, the wrong direction. A check that the mean weight is 1 finds this at once.}

\iqfollow{Why is that weight a martingale in time?}{It is the exponential martingale $e^{\mu W_t - \mu^2t/2}$, with mean 1, and this is what makes the new measure a probability measure.}

\iqwhy{Change of measure is often stated abstractly. This version shows what it does in one line, and gives a practical reason to care.}
\end{solution}

\begin{solution}{iq:iv1:stochastic-calculus:18}
Real-world $0.655$, risk-neutral $0.460$, price $0.460$. The market price of risk is $\frac12$. The price does not depend on $\mu$.

$\ln S_1 - \ln 100$ is normal with standard deviation $0.2$. Its mean is $\mu - \frac12\sigma^2 = 0.08$ in the real world, so $P(S_1 > 100) = \Phi(0.08/0.2) = \Phi(0.4)$.

Under the risk-neutral measure, the drift is the rate, here zero\index{risk-neutral measure}, so the log mean is $-0.02$, and the probability is $\Phi(-0.1) = 0.460$. With a zero rate, the price is this probability, by the replication argument for a digital.

The market price of risk is $\lambda = (\mu - r)/\sigma = 0.5$\index{market price of risk}. The two measures are related by $dQ/dP = e^{-\lambda W_1 - \lambda^2/2}$, which removes the drift $\lambda\sigma$ from the return. Reweighting real-world paths with this factor gives $0.460$ in a simulation, and the weight has mean 1.

Anyone holding the digital at $0.46$ has a positive expected profit, of $0.655 - 0.46$. This is not free: the digital pays in states where the stock is up, so it carries stock risk, for which the holder expects compensation.

\iqwrong{Pricing with the real-world probability, or thinking that $\mu$ affects the price because it affects the chance of the payoff. Hedging removes it from the price, though it does not remove it from the holder's expected profit.}

\iqfollow{What are the probabilities if $\mu = 2\%$?}{The real-world chance is now $\Phi(0) = \frac12$, and the price is still $0.460$. The gap has shrunk to $0.04$, because the expected profit comes only from the drift.}

\iqwhy{Confusion between probabilities and prices is the most common error in this area, and the digital isolates it.}
\end{solution}

\begin{solution}{iq:iv1:stochastic-calculus:19}
$u(t, x) = e^{-r(T - t)}\mathbb{E}[g(X_T) \mid X_t = x]$, where $dX = \mu\,dt + \sigma\,dW$. The two solutions are $x^2 + T - t$ and $e^{x + (T - t)/2}$.

Let $M_s = e^{-rs}u(s, X_s)$. Ito's lemma gives drift $e^{-rs}(u_t + \mu u_x + \frac12\sigma^2u_{xx} - ru)$, which is 0 by the equation\index{Feynman-Kac formula}. So $M$ is a martingale, and $M_t = \mathbb{E}[M_T \mid X_t]$, which is the formula.

With $\mu = r = 0$, $\sigma = 1$: given $W_t = x$, $W_T$ is normal, with mean $x$ and variance $T - t$.

For $g = x^2$: $u = x^2 + (T - t)$. Check: $u_t = -1$ and $\frac12u_{xx} = 1$, which sum to 0.

For $g = e^x$: $u = e^{x + (T - t)/2}$. Check: $u_t = -\frac12u$ and $\frac12u_{xx} = \frac12u$, which also sum to 0.

Finite differences and sampling agree with both at the points tried. A rate $r$ multiplies each by $e^{-r(T - t)}$, and this can be checked by substitution. The terminal value is $g(x)$, as required.

\iqwrong{Getting the sign of $u_t$ wrong. This is a backward equation, solved from the terminal time. So time to go, $T - t$, appears, and not $t$.}

\iqfollow{How do you get the Black--Scholes equation from this?}{Use $dS = rS\,dt + \sigma S\,dW$ under the risk-neutral measure. Then the formula says that the price solves $u_t + rSu_S + \frac12\sigma^2S^2u_{SS} - ru = 0$, which is that equation.}

\iqwhy{This is the link between expectations and differential equations, which lets you choose the easier of the two.}
\end{solution}

\begin{solution}{iq:iv1:stochastic-calculus:20}
$dX = X\,dW$. The mean is always 1, the median is $e^{-T/2}$, and $P(X_T < 1) = \Phi(\sqrt{T}/2)$, which is $0.69$ at $T = 1$ and $0.9999997$ at $T = 100$. Almost every path goes to 0.

$X$ is the exponential martingale, with $dX = X\,dW$ by Ito's lemma. So the mean stays at 1.

$\ln X_T = W_T - T/2$ is normal with mean $-T/2$ and variance $T$. So the median is $e^{-T/2}$: $0.607$ at $T = 1$ and about $2 \times 10^{-22}$ at $T = 100$. Also, $X_T < 1$ needs $W_T < T/2$, whose chance is $\Phi(\sqrt{T}/2)$\index{exponential martingale}.

Since $W_t/t \to 0$, $\ln X_t \approx -t/2 \to -\infty$, so each path goes to 0 with probability one.

The reconciliation is that the mean is carried by ever rarer paths with enormous values. The variance is $e^T - 1$, which is huge. At $T = 100$, no feasible simulation can show a mean of 1. A million paths usually give far less, with the rest in a tail too thin to sample. Note that optional stopping does not apply at $T = \infty$: the limit is 0, while the mean is 1, so the martingale is not uniformly integrable.

In investment terms, this is a fair bet at full size, with a positive chance of huge gains and almost sure ruin in the long run. That is the reason for caring about growth rates and not only expected values.

\iqwrong{Saying that the mean is a good guide to typical wealth. It is the extreme example of a mean that describes no path anyone is likely to see.}

\iqfollow{What are the second moments at $T = 1$ and $T = 100$?}{$e^1 = 2.72$ and $e^{100} = 2.7 \times 10^{43}$. So the sample mean is useless long before $T = 100$.}

\iqwhy{This separates ensemble averages from time averages, which is the heart of the discussion of sizing and of the Kelly criterion.}
\end{solution}

\begin{solution}{iq:iv1:stochastic-calculus:21}
The conditional mean is $TW_T/2$, and the conditional variance is $T^3/12$. At $T = 1$, these are $W_1/2$ and $\frac{1}{12}$, and the endpoint explains $\frac34$ of the variance.

$I_T$ and $W_T$ are jointly normal with mean zero, with $\operatorname{Var}I_T = T^3/3$, $\operatorname{Var}W_T = T$ and $\operatorname{Cov} = T^2/2$. These follow from $I_T = \int_0^T(T - t)\,dW_t$.

So normal regression\index{conditional normal} applies. The slope is $\operatorname{Cov}/\operatorname{Var}W_T = T/2$, which gives the mean $TW_T/2$. That makes sense: given the endpoint, the bridge is on average a straight line from 0 to $W_T$, whose integral is the area of a triangle.

The remaining variance is
\[
\frac{T^3}{3} - \frac{(T^2/2)^2}{T} = \frac{T^3}{3} - \frac{T^3}{4} = \frac{T^3}{12}.
\]
A second route: it should equal the variance of the integral of a Brownian bridge on $[0, T]$. For $T = 1$, $\int_0^1\int_0^1(\min(s, t) - st)\,ds\,dt = \frac13 - \frac14 = \frac{1}{12}$, the same.

The share explained is $(\frac13 - \frac{1}{12})/\frac13 = \frac34$, which is $0.866^2$, the squared correlation between the two. A simulation with $40\,000$ paths recovers the slope, the residual variance and the share.

\iqwrong{Thinking that the integral is independent of the endpoint because the increments are independent. They are not independent: both are sums of the same increments.}

\iqfollow{Does this give a control variate for an Asian option?}{Yes. Since the endpoint explains three quarters of the variance, a simulated average can be regressed on $W_T$, and that alone cuts the variance by that amount. A geometric average does better still.}

\iqwhy{Conditioning on a normal variable is a one-line tool, and the same tool appears in Brownian bridge constructions and in variance reduction.}
\end{solution}

\begin{solution}{iq:iv1:stochastic-calculus:22}
$A$ has the arcsine law, with density $1/(\pi\sqrt{x(1 - x)})$: mean $\frac12$, variance $\frac18$, and $P(A \le 0.1) = \frac{2}{\pi}\arcsin\sqrt{0.1} = 0.205$. So a skill-free trader is ahead for 90\% of the year about 20\% of the time.

The density is highest near 0 and 1, and lowest at $\frac12$\index{arcsine law}. The mean follows by symmetry. The variance comes from $\mathbb{E}A^2 = \frac38$, a double integral of the orthant probability.

The distribution function is $F(x) = \frac{2}{\pi}\arcsin\sqrt{x}$, so $P(A \le 0.1) = 0.205$. By symmetry, $P(A \ge 0.9)$ is the same. So with $41\%$ chance, the path spends more than 90\% of the time on one side. Compare a fraction between $0.45$ and $0.55$, which has a chance of only $6.4\%$. The most balanced outcome is the least likely.

A simulated walk with 2000 steps gives a mean of $0.50$, a variance of $0.125$ and a tail near $0.205$, which is a check independent of the formula. Note the conclusion. A profit curve above water for 90\% of the year says little about skill, as about one trader in five with no skill at all will have one.

A trader's edge, in contrast, moves the fraction reliably toward 1, but this measure is a weak test of it. A $t$-statistic on the mean uses the same data better.

\iqwrong{Assuming that a zero-mean path should spend about half its time above 0. The opposite is true: long stays on one side are normal, since excursions are long.}

\iqfollow{What is the standard deviation of $A$?}{$\sqrt{1/8} = 0.354$. That is above the uniform distribution's $0.289$, and a good way toward the $0.5$ of an equal mix of 0 and 1.}

\iqwhy{It is the strongest intuition-breaking result about random walks, with a direct reading for performance evaluation.}
\end{solution}

\begin{solution}{iq:iv1:stochastic-calculus:23}
$P(M_T \ge a, W_T \le x) = P(W_T \ge 2a - x)$. The gap $M_T - W_T$ has the law of $|W_T|$. The expected gap is 16\%, and the chance that it exceeds 20\% is $0.317$.

For $x \le a$, reflect each path after it first reaches $a$\index{reflection principle}. A path with $M_T \ge a$ and an end at or below $x$ corresponds one-to-one to a path ending at or above $2a - x$, which must have passed $a$ on the way. So the two events have equal probability.

Differentiating with respect to $a$ and $x$ gives the joint density of $(M_T, W_T)$. From it, the density of $M_T - W_T$ can be computed directly, and it is that of $|W_T|$.

A cleaner argument uses time reversal\index{time reversal}. Define $\hat{W}_s = W_{T - s} - W_T$, which is also a Brownian motion. Then $M_T - W_T = \max_s(W_s - W_T) = \max_s\hat{W}_{T - s} = \max_s\hat{W}_s$, whose law is that of the maximum, and so of $|W_T|$ by reflection. Hence
\[
\mathbb{E}(M_T - W_T) = \sqrt{2T/\pi}, \qquad P(M_T - W_T \ge c) = 2(1 - \Phi(c/\sqrt{T})).
\]

With volatility $\sigma$: multiply by $\sigma$. So $0.798 \times 0.2 = 16\%$, and $P(\text{gap} > 20\%) = 2(1 - \Phi(1)) = 0.317$. A simulation with $40\,000$ paths agrees after correcting for discrete checking. Since $\mathbb{E}W_T = 0$, the same number is also $\mathbb{E}M_T$.

\iqwrong{Treating a 16\% distance from the high as a signal of something wrong. In a driftless process with a 20\% volatility, it is what the average year looks like.}

\iqfollow{What is the expected gap after two years?}{It grows with the root of time: $16\% \times \sqrt{2} = 22.6\%$. Note that this is the current gap at the end. The maximum drawdown during the period is larger.}

\iqwhy{Time reversal is a beautiful trick, and it gives a benchmark for drawdowns that investors often misjudge.}
\end{solution}

\begin{solution}{iq:iv1:stochastic-calculus:24}
$\mathbb{E}\tau = a/\mu = 10$, and $\operatorname{Var}\tau = a\sigma^2/\mu^3 = 40$. $\tau$ has the inverse Gaussian law, with $P(\tau > 20) = 0.073$.

$\tau$ has finite mean under $\mu > 0$. For the mean, $X_t - \mu t$ is a martingale. Optional stopping\index{optional stopping} gives $a - \mu\mathbb{E}\tau = 0$.

For the variance, $(X_t - \mu t)^2 - \sigma^2t$ is a martingale. At $\tau$, $X_\tau = a$, so
\[
\mathbb{E}(a - \mu\tau)^2 = \sigma^2\mathbb{E}\tau.
\]
The left side is $\mu^2\operatorname{Var}\tau$, since $\mathbb{E}\mu\tau = a$. So $\operatorname{Var}\tau = \sigma^2(a/\mu)/\mu^2 = a\sigma^2/\mu^3$.

At the stated values: $\mathbb{E}\tau = 1/0.1 = 10$, and $\operatorname{Var}\tau = 0.04/0.001 = 40$, with a standard deviation of $6.3$.

This is the inverse Gaussian\index{inverse Gaussian distribution}, with mean $a/\mu$ and shape $a^2/\sigma^2 = 25$. Its tail has the closed form
\[
P(\tau \le t) = \Phi\left(\frac{\mu t - a}{\sigma\sqrt{t}}\right) + e^{2\mu a/\sigma^2}\Phi\left(\frac{-\mu t - a}{\sigma\sqrt{t}}\right),
\]
which gives $0.073$ for $P(\tau > 20)$. Sampling the law and a direct Euler run of the process agree with the mean and variance.

Limits: as $\mu \to 0$, both blow up, and the variance faster, as they must, since the driftless case has an infinite mean. At $\mu = 0.01$, the variance is already $40\,000$.

\iqwrong{Giving the deterministic time $a/\mu$ with no spread, or a variance of $\sigma^2/\mu^2$, which has the wrong scaling in $a$ and $\mu$. Dimensions catch that one.}

\iqfollow{If the target is 10 times as far, what happens to the relative spread?}{The mean and the variance each grow by 10, so the coefficient of variation falls by $\sqrt{10}$. Long-horizon targets are more predictable, proportionally.}

\iqwhy{Time-to-target is common in execution and in strategy evaluation. The martingale derivation is faster than the density.}
\end{solution}

\begin{solution}{iq:iv1:stochastic-calculus:25}
With $\lambda = \theta^2/2$, the martingale $e^{\theta W_t - \theta^2t/2}$ gives the result. The limit $\lambda \to 0$ shows that $T_a$ is finite, and the infinite derivative shows that its mean is infinite. At $a = 1$, $\lambda = \frac12$, the value is $e^{-1} = 0.368$.

For $\theta > 0$, $M_t = e^{\theta W_t - \theta^2t/2}$ is a martingale\index{exponential martingale}. Stopped at $T_a$, it is bounded by $e^{\theta a}$, as $W \le a$ before $T_a$ and the time term only reduces it. On the event $T_a = \infty$, $M_t \to 0$, as $W_t$ stays below $a$ while $t$ grows. So bounded optional stopping gives
\[
1 = \mathbb{E}\left[e^{\theta a}e^{-\theta^2T_a/2}\mathbf{1}(T_a < \infty)\right].
\]
Putting $\lambda = \theta^2/2$, we have $\mathbb{E}[e^{-\lambda T_a}\mathbf{1}(T_a < \infty)] = e^{-a\sqrt{2\lambda}}$.

The positive root matters. With $\theta < 0$, $M$ is unbounded as $W \to -\infty$ before $T_a$, so optional stopping fails. The other root would give a transform above 1, which is impossible.

Let $\lambda \to 0$: $P(T_a < \infty) = 1$. The derivative of $e^{-a\sqrt{2\lambda}}$ at 0 is $-\infty$, because of the square root, and so $\mathbb{E}T_a = \infty$.

At $a = 1$, $\lambda = \frac12$: $e^{-1}$. Sampling from $T_1 = 1/Z^2$ gives $0.368$.

\iqwrong{Applying optional stopping with no check on boundedness. The same step would show that $\mathbb{E}W_{T_a} = 0$, which is false, since $W_{T_a} = a$.}

\iqfollow{What changes with drift $\mu > 0$?}{The same argument with a shifted exponent gives $\mathbb{E}e^{-\lambda\tau} = e^{a(\mu - \sqrt{\mu^2 + 2\lambda})}$, with $\sigma = 1$. Now the derivative at 0 is finite, $-a/\mu$. At $\mu = \frac12$, $a = 1$, $\lambda = \frac12$, this gives $e^{0.5 - 1.118}$, as sampling confirms.}

\iqwhy{It is a compact use of the stopping theorem, and the choice of the root shows whether its conditions are understood.}
\end{solution}

\begin{solution}{iq:iv1:stochastic-calculus:26}
At $\mu = 10\%$, the expected time is $\ln 2/0.08 = 8.66$ years. At $\mu = 1\%$, the probability is $2^{-1/2} = 0.707$. At $\mu = 2\%$, it reaches 2 with certainty, but the expected time is infinite.

$X = \ln S$ has drift $\nu = \mu - \frac12\sigma^2$ and volatility $\sigma$, and the target is $a = \ln 2$\index{first passage time}.

$\mu = 10\%$: $\nu = 8\%$. The mean of a positive-drift hitting time is $a/\nu = 0.693/0.08 = 8.66$.

$\mu = 1\%$: $\nu = -1\%$, so the log drifts away. The chance of ever reaching a level $a$ above is $e^{2\nu a/\sigma^2} = e^{-0.5\ln 2} = 2^{-1/2}$. Note that the expected value of $S$ still grows.

Check with a martingale: $p = 1 - 2\mu/\sigma^2 = \frac12$, so $\sqrt{S}$ is a martingale\index{scale function}. It is bounded by $\sqrt{2}$ up to the hitting time, and if 2 is never reached, $S \to 0$ since $\nu < 0$, and the martingale tends to 0. So $1 = q\sqrt{2}$, which gives $q = 1/\sqrt{2}$. The answers agree, with no reference to the first formula.

$\mu = 2\%$: $\nu = 0$. The log is driftless, so it hits with certainty and has an infinite mean, as a one-sided hit for a driftless Brownian motion does.

A simulated 400-year run at $\mu = 1\%$ reaches 2 in about $68\%$ of paths after a correction for discrete checks. That matches the first-passage formula at that horizon, $0.683$, and the limit is $0.707$.

\iqwrong{Saying that any positive $\mu$ means the asset surely doubles. The sign that matters is that of the log drift. With $0 < \mu < \sigma^2/2$, the expected value rises while the typical path falls.}

\iqfollow{At which $\mu$ does the expected time become finite?}{For any $\mu > \sigma^2/2 = 2\%$, and it blows up at that value. The expected time at $\mu = 2.1\%$ would be about 693 years.}

\iqwhy{It tests the difference between arithmetic and geometric drift, and the check by a power martingale is a good habit.}
\end{solution}

\begin{solution}{iq:iv1:stochastic-calculus:27}
$|W_t| = \int_0^t\operatorname{sign}(W_s)\,dW_s + L_t$, with $\mathbb{E}L_1 = \sqrt{2/\pi} = 0.798$. $L$ measures time spent at 0, per unit of space, and it has the law of the running maximum.

A smooth approximation of $|x|$ has a second derivative that is concentrated around 0, with total mass 2. So the Ito term $\frac12f^{\prime\prime}(W)\,dt$ does not vanish, and it accumulates only when $W$ is at 0\index{Tanaka's formula}. Its limit is $L_t$, defined by
\[
L_t = \lim_{\varepsilon \to 0}\frac{1}{2\varepsilon}\int_0^t\mathbf{1}(|W_s| < \varepsilon)\,ds.
\]
This is the local time\index{local time} at 0: a rate of time spent near 0 per unit width.

The stochastic integral has mean zero, so $\mathbb{E}L_1 = \mathbb{E}|W_1| = \sqrt{2/\pi}$.

There are two estimators from a path. One is the occupation formula, with a small $\varepsilon$. The other is $|W_1| - \sum\operatorname{sign}(W_{i})\Delta W_i$, which rearranges Tanaka's formula. In a run, each has a mean near $0.798$, and the two estimators are strongly correlated across paths, above $0.9$. So they are measuring the same thing.

By Tanaka, $\int\operatorname{sign}(W)\,dW$ is itself a Brownian motion $B$, by L\'evy's characterisation, as its quadratic variation is $t$. Writing $|W_t| = B_t + L_t$ and noting that $L$ grows only when $|W| = 0$, Skorokhod's lemma gives $L_t = \max_{s \le t}(-B_s)$. So $L_t$ has the law of the maximum of a Brownian motion, and so of $|W_t|$. The quantiles from the run match.

\iqwrong{Applying Ito's lemma to $|x|$ with $f^{\prime\prime} = 0$, which gives $|W_t| = \int\operatorname{sign}(W)\,dW$. Its mean would be zero, which is impossible for a positive variable.}

\iqfollow{Where does this appear in hedging?}{A stop-loss strategy holds a share only when the price is above the strike. Its failure to replicate a call is exactly the local time at the strike: a cost that does not go away with more frequent trading.}

\iqwhy{This is one of the few places where the second-derivative term of Ito's lemma is visible on its own. It tests experienced candidates, and the hedging link is a real one.}
\end{solution}

\begin{solution}{iq:iv1:stochastic-calculus:28}
$\mathbb{E}Q = \frac12$, and $\operatorname{Var}Q = \frac13$. The series gives the same results, using also $\sum_k((k - \frac12)\pi)^{-4} = \frac16$.

The mean is $\int_0^1\mathbb{E}W_t^2\,dt = \int_0^1t\,dt = \frac12$.

For the variance, $\operatorname{Cov}(W_s^2, W_t^2) = 2\min(s, t)^2$, as for any two normals with mean zero. So
\[
\operatorname{Var}Q = 2\int_0^1\int_0^1\min(s, t)^2\,ds\,dt = 4\int_0^1\int_0^ts^2\,ds\,dt = 4\int_0^1\frac{t^3}{3}\,dt = \frac13.
\]

For the series, the sine functions $\sqrt{2}\sin((k - \frac12)\pi t)$ are orthonormal on $[0, 1]$\index{Karhunen-Loeve expansion}. Writing $\lambda_k = ((k - \frac12)\pi)^{-2}$ and integrating the square of the expansion,
\[
Q = \sum_k\lambda_kZ_k^2,
\]
since cross-terms vanish. So $\mathbb{E}Q = \sum_k\lambda_k = \frac12$, as given, and $\operatorname{Var}Q = 2\sum_k\lambda_k^2$, as $\operatorname{Var}Z^2 = 2$. Thus the required sum is $\sum_k\lambda_k^2 = \frac16$. This is the classical sum $\sum 1/(2k - 1)^4 = \pi^4/96$ in disguise, and a numeric sum gives $\frac16$ to twelve decimals. So the two routes agree.

The first term has $\lambda_1 = 4/\pi^2 = 0.405$, and it carries $98.6\%$ of the variance. So $Q$ is close to a scaled chi-square with one degree of freedom, and its distribution is strongly skewed. A simulation returns $0.50$ and $0.33$.

\iqwrong{Using $\mathbb{E}Q^2 = (\mathbb{E}Q)^2 + \int\operatorname{Var}W_t^2\,dt$, which ignores the correlation over time. It gives $\frac14 + \frac23$ in place of $\frac14 + \frac13$, which is too large.}

\iqfollow{Why does a single term carry so much of the variance?}{The eigenvalues fall as $1/k^2$, and their squares as $1/k^4$. A random walk's level is dominated by its slowest component.}

\iqwhy{Such quantities arise in unit-root statistics, whose limits are integrals like $Q$. The expansion tells you much more than the two moments do.}
\end{solution}

\begin{solution}{iq:iv1:stochastic-calculus:29}
$u(0, 0) = 1/\sqrt{\cosh\theta}$, with $A = \theta\tanh(\theta(1 - t))$, and $B$ its integral. At $\theta = 1$, the value is $0.805$.

Define $u(t, x) = \mathbb{E}[e^{-\frac12\theta^2\int_t^1W_s^2\,ds} \mid W_t = x]$. Feynman--Kac\index{Feynman-Kac formula} with a state-dependent discount gives
\[
u_t + \frac12u_{xx} - \frac12\theta^2x^2u = 0, \qquad u(1, x) = 1.
\]
Substitute $u = \exp(-\frac12A(t)x^2 + B(t))$. Then $u_t = (-\frac12A'x^2 + B')u$ and $u_{xx} = (A^2x^2 - A)u$. Matching powers of $x$:
\[
A' = A^2 - \theta^2, \qquad B' = \tfrac12A, \qquad A(1) = B(1) = 0.
\]
The Riccati equation\index{Riccati equation} has the solution $A(t) = \theta\tanh(\theta(1 - t))$. Check: $A' = -\theta^2\operatorname{sech}^2 = \theta^2(\tanh^2 - 1) = A^2 - \theta^2$.

Then
\[
B(0) = -\int_0^1\frac12\theta\tanh(\theta(1 - t))\,dt = -\frac12\ln\cosh\theta.
\]
So $u(0, 0) = e^{B(0)} = (\cosh\theta)^{-1/2}$\index{Cameron-Martin formula}.

For small $\theta$, this is $1 - \theta^2/4$, while the left side expands to $1 - \frac12\theta^2\mathbb{E}Q = 1 - \theta^2/4$. They agree. As $\theta \to \infty$, it decays like $\sqrt{2}e^{-\theta/2}$, so the chance of a very small $Q$ falls off quickly, as it must.

A second route uses the series: $Q = \sum\lambda_kZ_k^2$ gives $\prod_k(1 + \theta^2\lambda_k)^{-1/2}$, and the product form of the hyperbolic cosine turns this into the same result. A numerical product and a simulation both give $0.805$.

\iqwrong{Trying a guess with no $x^2$ term, or getting the sign of the $\tanh$ wrong so that $u$ grows with $x$. The discount punishes large $x$, so $u$ must fall with it.}

\iqfollow{Why is a quadratic exponent the right guess?}{The potential is quadratic and the diffusion is Gaussian. Such problems stay Gaussian, and a quadratic exponent is the most general Gaussian form.}

\iqwhy{This requires several techniques at once. The Riccati structure is the same one behind quadratic-Gaussian and affine models.}
\end{solution}

\begin{solution}{iq:iv1:stochastic-calculus:30}
$\mathbb{E}e^{-\lambda\tau} = 1/\cosh(a\sqrt{2\lambda})$, $\mathbb{E}\tau = a^2$, $\operatorname{Var}\tau = \frac23a^4$, and $P(\tau > t)$ decays at rate $\pi^2/(8a^2)$.

$\cosh(\theta W_t)e^{-\theta^2t/2}$ is a martingale, as the average of $e^{\pm\theta W_t - \theta^2t/2}$\index{exit time}. It is bounded up to $\tau$, and at $\tau$ we have $|W_\tau| = a$, so $\cosh(\theta W_\tau) = \cosh\theta a$. Optional stopping gives $\cosh(\theta a)\mathbb{E}e^{-\theta^2\tau/2} = 1$, and with $\lambda = \theta^2/2$ this is the result.

Expanding with $x = a\sqrt{2\lambda}$: $1/\cosh x = 1 - \frac12x^2 + \frac{5}{24}x^4 - \dots$, and $x^2 = 2\lambda a^2$. So
\[
\mathbb{E}e^{-\lambda\tau} = 1 - \lambda a^2 + \frac56\lambda^2a^4 - \dots
\]
Matching with $1 - \lambda\mathbb{E}\tau + \frac12\lambda^2\mathbb{E}\tau^2$: $\mathbb{E}\tau = a^2$, and $\mathbb{E}\tau^2 = \frac53a^4$. So $\operatorname{Var}\tau = \frac23a^4$.

Check: the polynomial is a martingale\index{Hermite martingales}. At $\tau$, $a^4 - 6a^2\mathbb{E}\tau + 3\mathbb{E}\tau^2 = 0$. With $\mathbb{E}\tau = a^2$, this gives $\mathbb{E}\tau^2 = \frac53a^4$, the same.

For the tail, $\cosh(a\sqrt{2\lambda})$ vanishes when $a\sqrt{2\lambda} = i\pi/2$, so at $\lambda = -\pi^2/(8a^2)$. The nearest singularity sets the exponential decay of the tail\index{tail decay}: $P(\tau > t) \approx \frac{4}{\pi}e^{-\pi^2t/(8a^2)}$.

At $a = 1$, $\lambda = \frac12$: $1/\cosh 1 = 0.648$, and $P(\tau > 4) \approx 0.0092$. A symmetric lattice walk with $100\,000$ paths agrees on the mean, the variance, the transform and the tail. Compare with the one-sided case: a second barrier turns the heavy tail, of order $1/\sqrt{t}$, into an exponential one, and every moment becomes finite.

\iqwrong{Using $\mathbb{E}\tau^2 = (\mathbb{E}\tau)^2 + \mathbb{E}\tau$, or dropping the factor $\frac12$ in the expansion of the exponential. The polynomial martingale is a quick check on this kind of slip.}

\iqfollow{Why is $\frac{4}{\pi}$ the constant in the tail?}{The full tail is the series $\frac{4}{\pi}\sum_{k \ge 0}\frac{(-1)^k}{2k + 1}e^{-(2k + 1)^2\pi^2t/(8a^2)}$, and the first term dominates quickly. At $t = 0$ the alternating sum is $\pi/4$, which gives 1, as it must.}

\iqwhy{It connects a martingale, a transform, moments and a tail in one problem. The check by a second polynomial martingale shows the habit expected at the top level.}
\end{solution}
